\subsection{Conclusions}

Inferring population and diversification dynamic processes from biological data is essential for understanding the forces shaping biodiversity and the spread of pathogens across space and time.
This task is inherently challenging because available data are incomplete: only a small fraction of ancient taxa leave a fossil record, many extant species remain unsampled, and in epidemic settings only a subset of infected individuals is typically sequenced.
Phylodynamic methods provide a principled framework for reconstructing past dynamics from such incomplete data.
However, as models become more complex, parameter uncertainty can increase to the point that inferences lose interpretability, and mechanistic links between key processes---such as speciation, extinction, and migration---and environmental or biotic drivers remain difficult to quantify.
Our BELLA approach enables accurate, data-driven reconstruction of complex phylodynamic patterns while controlling overfitting and retaining interpretability through explainable AI techniques.

We envision this framework enabling robust phylodynamic inference across a wide range of application domains, including macroevolution and epidemiology, as demonstrated here, as well as areas such as single-cell biology and linguistics where phylodynamic tools are increasingly applied.
Applications naturally include phylogeography, which in our formulation falls within the broader phylodynamic framework.
More generally, this approach opens the door to incorporating mechanistic relationships between polulation dynamic parameters and diverse predictors, ranging from life-history traits to environmental variables.
Together, these developments provide a foundation for more general and mechanistically grounded phylodynamic inference, advancing our understanding of the processes that shape population dynamics and evolution.